\medskip
\noindent\textbf{Additive basis of order $r$.}
A set $A\subseteq\mathbb{N}$ is a \emph{basis of order $r$} if every sufficiently large integer $n$ can be written as a sum of at most $r$ (not necessarily distinct) elements of $A$.

\medskip
\noindent\textbf{Exact order.}
A set $A$ has \emph{exact order} $k$ if every sufficiently large integer is representable as a sum of exactly $k$ (not necessarily distinct) elements of $A$.

\medskip
\noindent\textbf{Extremal function.}
For each $r\ge 2$, define $h(r)$ as the largest finite $k$ such that there exists a basis $A$ of order $r$ that has exact order $k$.
The problem asks for the value of
\[
\lim_{r\to\infty} \frac{h(r)}{r^2},
\]
and for the exact values of $h(r)$.
